% Figure: conduction-band staircase of a quantum-cascade laser under bias. The
% band tilts under the applied field; in each period an electron drops between
% two subbands of the same conduction band, emitting a photon, then tunnels to
% the next period and repeats. Schematic; all coordinates entered explicitly.
\begin{figure}[htbp]
\centering
\begin{tikzpicture}[font=\scriptsize, x=1cm, y=1cm]
% overall tilted conduction-band edge (applied field drops potential to the right)
\draw[enggray, thick] (0,4.2) -- (10,0.2);
\node[enggray, anchor=south, font=\tiny] at (2.0,3.6)
  {conduction-band edge under bias};
% period 1 upper and lower subbands (short horizontal levels)
\draw[engblue, very thick] (0.6,3.9) -- (1.9,3.9);
\node[engblue, anchor=south, font=\tiny] at (1.25,3.9) {upper};
\draw[engblue, very thick] (0.9,3.15) -- (2.2,3.15);
\node[engblue, anchor=north, font=\tiny] at (1.55,3.15) {lower};
% radiative transition (photon) within period 1
\draw[engred, ->, very thick] (1.5,3.85) -- (1.7,3.20);
\node[engred, anchor=west, font=\tiny] at (1.75,3.55) {$h\nu$};
% tunnelling injection to next period
\draw[engorange, ->, thick, dashed] (2.2,3.15) -- (3.5,2.9);
% period 2
\draw[engblue, very thick] (3.5,2.9) -- (4.8,2.9);
\draw[engblue, very thick] (3.8,2.15) -- (5.1,2.15);
\draw[engred, ->, very thick] (4.4,2.85) -- (4.6,2.20);
\node[engred, anchor=west, font=\tiny] at (4.65,2.55) {$h\nu$};
\draw[engorange, ->, thick, dashed] (5.1,2.15) -- (6.4,1.9);
% period 3
\draw[engblue, very thick] (6.4,1.9) -- (7.7,1.9);
\draw[engblue, very thick] (6.7,1.15) -- (8.0,1.15);
\draw[engred, ->, very thick] (7.3,1.85) -- (7.5,1.20);
\node[engred, anchor=west, font=\tiny] at (7.55,1.55) {$h\nu$};
\draw[engorange, ->, thick, dashed] (8.0,1.15) -- (9.3,0.9);
% one electron path label
\node[engorange, anchor=north, font=\tiny] at (5.0,0.55)
  {one electron emits one photon per period};
\draw[engorange, ->] (5.0,0.7) -- (5.6,1.6);
% axes hints
\draw[enggray, ->] (0,-0.1) -- (10,-0.1);
\node[enggray, anchor=north, font=\tiny] at (5,-0.15) {position (growth direction)};
\draw[enggray, ->] (-0.15,0) -- (-0.15,4.4);
\node[enggray, anchor=south, rotate=90, font=\tiny] at (-0.4,2.2) {energy};
\end{tikzpicture}
\caption[Quantum-cascade-laser conduction-band staircase]{Conduction-band
profile of a quantum-cascade laser under operating bias. The applied field tilts
the band into a staircase of nominally identical periods. Within each period an
injected electron sits in an upper subband and drops to a lower subband of the
same conduction band, emitting a photon whose energy is the subband spacing, set
by the well widths rather than by any material band gap. The electron then
tunnels into the upper subband of the next period and repeats, so a single
electron emits one photon per period and a stack of tens of periods multiplies
the gain. Because the transition is intersubband, the emission wavelength is
tuned by layer thickness across the mid- and far-infrared, where no ordinary
interband diode reaches.}
\label{fig:vol04ch12:qcl-staircase}
\end{figure}
